\chapter{边界层理论基础}
当粘性流体绕流的特征雷诺数很大时，在物体表面形成粘性起主要作用的薄层， 称作边界层.在边界层内流动必须考虑流体粘性;边界层外流动可以用无粘的欧拉方 程计算.边界层的发现对于真实流体绕流的分析和计算有极大的推动作用.本章介绍 边界层理论，以及比较简单的边界层和其他粘性薄层的分析与计算方法.

\section*{10.1  牛顿流体大雷诺数的定常绕流}\addcontentsline{toc}{section}{10.1 牛顿流体大雷诺数的定常绕流}

自然界和工程中许多流动问题的雷诺数很大，如飞行器周围的外部流动或流体 机械的内部流动都属大 Re 流动问题.举例来说 z 小型的低速飞机，飞行速度 U$\infty$ = 360km/h. 机身长度为 30m. 则在气温 293K 大气中飞行时 ($\nu$= 1. 5 X 10- sm2 /s): \textasciitilde{} 100 X 30 "唱，..:; n ... \textasciitilde{}唱 d 1'(e = 一$\tau$$\tau$一 $\lambda$lV ' = \textasciitilde{}$\lambda$lV - 同样速度和长度的高速水翼船在水中航行时，由于水的运动粘性系数较空气小一个 量级，约为 $\nu$= 1. 14X lO -- 6 m 2 /s. 所以雷诺数更高 z 100 X Re = 寸. 14一 X 10 6

\[= 2. 7 X 10\]

9 虽然大雷诺数流动意味着粘性应力相对于惯性力很小，但是完全忽略粘性而把 流体近似为理想流体时，就会产生绕流物体的阻力为零的矛盾 .20 世纪初著名流体 力学家普朗特首先发现，高雷诺数绕流时，粘性只在贴近物面的极薄一层内主宰流体

运动，并称这一薄层为附面层或边界层.在边界层以外，流动仍可视作理想的.这一发 现成为近代空气动力学基础，并在许多流体力学问题中得到广泛的应用，使经典的流 体力学成为实用的技术学科.作为纳维『斯托克斯方程的另一类近似解，可以用摄动 方法从理论上探讨高雷诺数的流动特征和求解方法. 1.高 Re 数流动常规摄动的奇异性 对大雷诺数流动，我们取 E; =l/Re 为小参数对流动变量作常规的摄动展开 z Ui = U$\tilde{I}$) 十 E; U\textasciitilde{}2) + O( E;2 )

\[\rho =\rho <))+\psi (2)\]

代入定常的不可压缩牛顿流体运动方程，可得 "、 J 盯$\tilde{I}$) /...T(?) aU$\tilde{I}$) I rT( 1) aU\textasciitilde{}2) 气，Ujl\}T$\iota$ +dUjZ) $\tau$一一 + U\textasciitilde{} 1) ..:- ' J +${\rm ο}$ ( E; 2 ) cJ Xj 飞 cJ Xj CIXj I =一纽工 -E; 纽工 +E; 主旦二十 O( E; 2 )

\[Jx, - JXi ,- JXjdXj\]

方程的一阶近似为 JU\textasciitilde{})) ， aU\textasciitilde{}2) ， '" ? --，一 +$\epsilon$ v\_ ， 一十 O( E;<) = 0 rJ Xi clXi 以))旦二一丝二一- J r1 Xj JXi

\[aU:1> \wedge \]

aXi 它们是熟知的欧拉方程组.也就是说:大 Re 流动的常规摄动近似是理想流动， 但是欧拉方程只能满足固壁不穿透条件，而不能满足固壁元滑移条件.所以，常规摄 动展开得到的低阶近似方程在固壁边界处有奇性，它不能描述壁面处的真实物理现 象.普朗特对此提出了"满足壁面附近粘附条件的粘流解与远离壁面无粘解的渐近衔 接方法"，即边界层理论. 2. 普朗特理论一一 有粘、无粘流动的渐近衔接方法 德国哥丁根学派的普朗特教授于 1904 年以铝粉为示踪剂进行了大 Re 绕流的流 动显示研究.他发现:贴近物面处流体运动速度很慢，在与物面很小的距离以外流线 谱与理想的位势理论预测结果基本一致.普朗特根据这一发现，提出了边界层理论概 念一一定常绕流申流体粘性只在贴近物面极薄的-层内主宰流体运动，称这-层为 边界层;边界层外的流动可近似为无粘的理想流动. 他进一步用量级分析方法导出了边界层方程，建立了绕流场，它可由满足欧拉方 程的外部流场与有粘性应力的边界层方程控制的内部流场组合而成.普朗特的这一 发现不仅为流体力学的应用开辟了正确的道路，而且也为含小参数非线性微分方程

求解提供了 一 种新方法.应用数学家们在普朗特量级分析的基础上发展了渐近衔接 摄动方法，这种方法实际上是普朗特思想的数学表达，为了便于叙述清楚这 一 理论， 以平壁附近的不可压缩流体的定常绕流为例来说明普朗特理论的思想和衔接方法， 并导出边界层方程. (1)两种尺度的分区流动现象 设有均匀定常来流绕过极薄的平板，流动的特征雷诺数很大.按照理想流体的无 量 I ，，\textasciitilde{} x 粘流动，极薄平板对均匀流动毫无影响，因为 流体可以在平艇上滑移过去.事实上，平板表 面的流体速度等于零，在远离平板处流动才恢 复到均匀流动.假定平板展向无限长，该流动 是 二 维的，如图10. 1 所尽. 图10.1 平壁面绕流的边界 层 普朗特理论的基本思想是在大 Re 绕流 中存在两个流动区域:外民是常规几何尺度 L( 例如平板的长度)和流动速度尺度 U 、 ;内区限于贴近固壁很小几何尺度(们的区域里，内区流向尺度为 L ，横向尺度为 ${\rm o}$-d- 也<< 1 ).根据这 一 思想，内外场用不同特征长度构成无量纲表达式. 外区用常规尺度作无量纲化(用上标"眷"表示无量纲量) : u = U . u ' (.r ' .y ' ). v = U . v ' (.r ' .y') ， $\rho$ = 严 $\rho$ . (.r . ， y 帽 ) 其中

\[I ' = .r/L. y ' = y / L\]

内区(边界层内，用上标"，\_\_，"表示无量纲量)的流向是常规几何尺度，横向是边界 层尺度 E L( E << l) .即流向和横向的无量纲坐标为 x = .r 暑 = .r/ 1-, j = y 嗣 / E = y / (d\_) 内区流场应 写 作: u = U .. ${\rm u}$( x . j ). v = U ${\rm ‘}$ 专 ( x .j) ${\rm •}$ $\rho$ = ,xP.. ${\rm ß}$ (:i , j) 内、外场采用两种横向尺度表明，同一几何坐标 y ， 在边界层内放大了 1 /$\epsilon$ 倍.对 于给定 y 值，有 j >> y. ${\rm •}$ 就是说同 - y 值对于外区来说它的无量纲坐标 y ' 很小，可 以认为非常贴近固壁;但是在边界层坐标中主值很大.形象地说，用放大镜来观察边 界层内的流动，而用普通平光镜来观察外区流动.为了使内外区近似解在全场连续过 渡.内外区流动应满足渐近衔接条件 z 内区解的外极限。$\rightarrow$$\infty$ )应和外区解的内极限 勺 ' $\rightarrow$ 0) 相等，即 i 可 (.r ' ， y') = l 川 j) limp 篝 (.r 幡 ， y' ) = lim${\rm ß}$( 牙 .j) .Y ... 0 y- 町 (10. 1) 以上的分区表达式和内外场渐近衔接式构成了边界层理论的框架.下面以 二 维寇常 流动为例导出边界层流动方程. 35 二

边界层理论基耐 (2) 外场和内场的摄动展开方程 (Re>>l) 将外场的无量纲式代入二维定常纳维-斯托克斯方程得 ${\rm •}$ ${\rm e}$1 u' I ${\rm •}$ ${\rm e}$J u' $\theta$$\rho$. 1 I ${\rm e}$1 2 U ${\rm •}$ I ${\rm O}$2 U' \textbackslash\{\} u -\_-\_;-刷 --\_一\textasciitilde{} --+--一 1-一一一;- ::--:;-;- I ${\rm O}$X 骨 I v ${\rm e}$1 y' ${\rm e}$1 X' I Re 飞 ${\rm e}$1 X'2 I iJy'2 J u. 纽二 +v' 主二=一丝二$\downarrow$土/主旦二 -$\iota$ 主旦二\}${\rm e}$1 X' I V ${\rm e}$J y' ${\rm e}$J y' I Re 飞 $\theta$ X. 2 ' ${\rm O}$y.2 J 00.2a) 第 10 章356 (10.2b) 些二+丝二 =0

\[<1X r1 Y\]

当 Re>>l 时，作为外场的近似式可略去动量方程的粘性项，即外场动量方程为欧拉 方程组.在均匀来流的情况下，外场是理想流体的定常有势流动:

\[a\phi a\phi \]

u = 一一一， v=-一 ${\rm e}$1 X' rJy ${\rm ‘}$ 00.3) $\rho$. = P( : 十七 1 一 1 V${\rm ø}$ 1 2 ) 将内场的元量纲表达式代人纳维-斯托克斯方程组得 $\iota$ ${\rm u}$ , v ${\rm e}$1 ${\rm u}$ rJ P I l' rJ 2 占 d'${\rm u}$ \textbackslash\{\} u 一+一一=一一 +$\tau$ $\leftarrow$\_\textasciitilde{} +寸二毛 l 3 主 E rJ ${\rm y}$. ${\rm e}$1:; l<e 飞rJ P E" ${\rm e}$1 ${\rm y}$' , hv 一寸 y H I-ZE + 卜u-fJ-J It---

\[l-b +\]

一 、y 、。一、。1-E 一一 切一句$\tilde{u}$-E + 而

\[-M\tilde{u}\]

${\rm e}$1 ${\rm u}$ . 1 rJ v \_ 一?一一- .. ${\rm e}$J:; E ${\rm e}$1 ${\rm y}$ 由连续方程可以导出 z 内层的法向速度分量必须为 E 量级，专-()(E) .否则 $\epsilon$ $\rightarrow$0 时心 的渐近解等于零.因为 E $\rightarrow$0 时，连续方程近似为 rJ vl rJ y=O. 而由壁面粘附条件:豆 =0 时，专 =0. 从而内场处处古 =0. 这显然不符合内层流动的实际情况，因此内层法向速 度的量级应为 00.4) v (${\rm I}$ .${\rm y}$) = E U. 、 v (${\rm I}$ .${\rm y}$) 将式(1 0. 的代入内层纳维-斯托克斯方程后，内层的运动方程应写作

\[eJ u , - iJ u eJ p I 1 ,e12 u I 1 e1 2 u  \]

a 一 +v 一=一一 +$\tau$ 卜$\tau$+$\tau$ 一-$\tau$i${\rm e}$1:; ${\rm e}$1 ${\rm y}$ ${\rm e}$1:; l<e 飞 iJ:; $\iota$ E" ${\rm oy}$" (${\rm u}$ 而 ${\rm e}$1 v \textbackslash\{\} 1 ${\rm e}$1 P I E' 内 1 内一 一 1= 一一一+.:$\leftarrow$一 +$\tau$ 一丁${\rm e}$1:i ${\rm e}$1${\rm y}$ I $\epsilon$ ${\rm e}$1 ${\rm y}$ l< e 飞 $\theta$:;2 E' a${\rm y}$

\[nu --\]

古一 、y

\[3n-\]

$\tilde{u}$ 一 $\tilde{r}$ 吨。-、 d 1 a2 ii \_\_ ${\rm e}$1 2 ii 请注意第一式右边的粘性项，当 $\epsilon$<<1 时一一一》一一就是说横向的粘性扩散远远'$\epsilon$ 2a${\rm y}$2//a:;2' a2 ii 大于流向的扩散，因此一一可以略去.假设 ReE 2 -o( 1). 可令 E 2 = l/ Re. 得 z 方向运${\rm e}$1 :;2

动方程的近似式为 卜u 一$\rightarrow$ y 吨d- 、d + 布 -hM一- B 一 $\tilde{y}$

\[qo-\]

吨。

\[\tilde{U} +\]

$\tilde{u}$ 一 $\tilde{r}$ 吨d-

\[-u\]

将 ReE 2 =1 代人第二式后得 I d:U d :U 、 2 页 1 .1 2 :u 1 d2 :u \textbackslash\{\} E2 1 ${\rm u}$ 二+古\_\_:\_ 1 =一工 +E 4 -7+今一-7 1 飞 ${\rm o}$ :i rij I Jj 飞 riP E- dj $\zeta$/ 即 J${\rm þ}$ F"'td ?' 一- O(E') rij 于是，当 Re>> 1 和 E 2 Re= 1 时，边界层内渐近展开的 一 阶近似方程为 \_ d${\rm U}$ , \_ri${\rm U}$ rJ ${\rm þ}$ , 02 ${\rm u}$ u -一寸-v 一- -一-一寸-\_一 ri:i dj o:i Jj2 rJ ${\rm þ}$ 。 cJ j

\[cJ ii o:U .....:.:..+--:-=o o:i oj\]

00.5) 00.6) 00.7) 式 00.6) 表明: ri${\rm þ}$ / $\theta$ 豆 =0 ，即在边界层内，压强只在流向变化，沿边界层的法向压强 为常数，积分式(1 0.6) ，并应用渐近衔接式 00. 1) 得 ${\rm þ}$( :i ,j) = ${\rm þ}$ Ci， $\infty$ ) = ${\rm Þ}$' (.r * ,0) 上式表明边界层内的压强等于外区无粘流内边界的压强. 由导出的边界层方程<1 0.5)- 方程(1 0.8) 可归纳普朗特边界层理论的主要结 论如下. ${\rm 1}$边界层内压强在垂直壁面方向不变，沿壁面方向的压强等于外部流场的当地 壁面压强. ${\rm 2}$流向的分子粘性扩散远小于法向扩散，可以忽略.在数学上，这一简化导致不 可压缩流体定常流动的边界层方程由椭圆型的定常纳维 - 斯托克斯方程退化为抛物 (1 0.8) 型偏微分方程. ${\rm 3}$当 Re >> 1 时，边界层横向尺度 E$\tilde{l}$/ .jRe， 即边界层的横向尺度与 Re 的平方 根成反比. (3) 绕流问题的内外区桐合求解 应用边界层理论，我们可以构造高 Re 粘性绕流问题的位势流和边界层流动的 组合解.具体步骤如下. ${\rm 1}$首先用理想流边界条件求解物体绕流的位势流方程，得到全场无粘流解 z u' (.r .y). 扩 (.r .y) 和 $\rho$ ${\rm ‘}$ (.r， y) ， 这一步骤的具体解法已在第 5 章讲述. ${\rm 2}$以无粘流解物面的流动参数为边界层方程的外边界条件，求解边界层方程.

具体来说，将无粘流边界的压强按伯努利公式算出: p' (.r, 0 ) = p,: - u; 2 (.r, 0) /2 (式中 : u; =U./U山 ， U. 是外部无粘绕流在物面的速度) ，它应等于边界层内的压强 pCi) = p' (.r, 0) =到 - u; 2 (x , 0) /2 ，将其代入边界层方程，得 \_ rJ ${\rm i}$i 1 、 $\theta$ 证 ${\rm •}$ du; 1 ${\rm e}$J 2 ${\rm i}$i u- 寸-v - = u\_ -.-寸-- rJ:i $\theta$ 豆. d.r ' ${\rm e}$J y2

\[eJ ii + Jv = 0\]

${\rm e}$J:i ${\rm e}$J y 边界层方程的边界条件为 z 壁面无滑移条件

\[y = O:ii = V = 0\]

内外区衔接条件

\[y- \infty  ,  ii = u; (.r)\]

将以上边界层方程写回到有量纲形式为 边界条件是 i) u 月 u ${\rm ••}$ dU\_ , ${\rm e}$J2 U u ::" + v 二 =U 一:..!.+$\nu$ 一号${\rm e}$J.r' -()y -, d.r '\textasciitilde{}${\rm e}$J y2 ${\rm e}$J u , ${\rm e}$) 啊 士主 +$\tau$二 =0

\begin{gather*}
<I.r r1 y\\
y=O , U=V=O\\
y \rightarrow \infty  ,  u = U,(.r)
\end{gather*}

3. 边界层厚度的各种定义 (1 0.9) (1 0.10) (1 0.11) (1 0.12) 从边界层方程的建立过程，可看出边界层方程的解有这样的性质 z 当 y v但"${\rm e}$/L$\rightarrow$$\infty$时，流向速度 u 渐近地达到外流速度 U.. 实际上 ， u $\rightarrow$U， 的渐近过程非 常迅速，当 y 大于某一值时 .u 与 U， 的差别微乎其微，粘性的影响就可以忽略.因此 常常以流向速度渐近的程度来定义边界层厚度 8. 通常沿壁面的法线向外推移，以 u=0.99U， 位置和壁面间的距离定义为边界层厚度，并称为边界层的名义厚度，如 图 10.2(a) 所示. v A 11 (a) (b) 图 10.2 边界层厚度和位移厚度示意阁 (a 】名义厚度和位移厚度; (b) 位移厚度的-种解释

用渐近性质来定义的边界层名义厚度，在实验测量或理论计算中常常会有较大 计算误差.为了准确起见，采用另外一些方法来定义边界层的厚度.下面给出另外两 种边界层厚度的定义. 定义 10.1 边界层的排挤厚度.用下式定义边界层的排挤厚度(也称位移厚 度 )8 1 : 81=jf(1-u川dy 00.13) 定义 10.2 边界层的动量损失厚度.用下式定义边界层的动量损失厚度 82 : $\delta$2=jutsV/Ue)(I-u 川 dy (\}O$\omega$ 显然，在边界层任意流向位置上，以上二式右边的职分都有确定的数值.由于边 界层内速度分布的渐近性，上限取为 8 或者$\infty$，积分值的误差很小，所以积分上限括 号中写 一 个 8. 表示也可用 $\delta$ 为积分上限.下面讨论乱和 82 的意义. (1)排挤厚度 81 的物理意义 排挤厚度(义称位移厚度 )8 1 的计算式为 Sl=jJm[ l-u 川.2') Jdy 考察边界层流动的速度分布，如图 10.2(a) 所示.在壁面上速度为零，在边界层外边界 流体速度达到外场势流速度 U，. 由于壁面附近流体速度的减小，理想均匀来流(图 10.2 (a) 中的 ABcm 进入到边界层后.遵循不可斥缩流体定常流动体积流量相等的规律. 它的外边界将向外偏移(图 10.2(a) 中的流线 CC'). 形象地说，由于流体的粘性而形 成的边界层将理想流动的流线向外排挤，或者说理想流动的流线向外位移.可以根据 流量连续的原则来计算理想流动流线向外的位移量在边界层内，它的流量是 r udy (图10.2(a) t!'面积 A'C'D'). 从边界层上游理想位势流动的来流是均匀的速度 U，. 这时要通过与边界层内同样多的流量(图 10.2(a) 中面积 ABCD) 理想流动所需横 向厚度为 8- 1::l. 8. 并定义横向距离1::l.8 为位移厚度 81 ，由不可压缩流体的连续性方 程得 即 f: 向 =ftdy=jJudy-ASU， AS=SI =-i一 r'

\begin{gather*}
'8>\\
[U,(.r) - uJdy U, Lr)J " L- '''-- --,---,\\
=J.<SUJr)) 11-L]
\end{gather*}

以上推导说明 8 1 的物理意义:厚度为 (8 一 $\delta$1 )的理想位势流进入边界层后，由于 近壁流速减小，它的外边界外移，相当于物面增加厚度仇，故 8 1 称为位移厚度或排挤

厚度根据位移厚度公式，位移厚度还可以作另一种解释 [(6】肌 x) 一山表示 由于流体粘性，厚度为 S 的理想流体流入边界层后所损失的流量，如果这股流量以理 想流体的均匀速度流动，它的厚度等于位移厚度.图 10.2(b) 直观地说明位移厚度的 几何意义图中画右斜线部分的面积等于 L 咄 [U， $\omega$ 一巾，它等于同左斜线的短 形面积 U.${\rm ð}$ l. 从这个意义上来说，位移厚度也可称作边界层的质量损失厚度. (2) 动量损失厚度 82 的物理意义 动量损失厚度 82 的计算公式是

\[82 = r\]

W 旦 11- 旦 ldv < J 0 u， 飞 u， r J 现在考察边界层内的动量通量以及它和边界层上游流量相同的位势流动量通量之 差.边界层内流体的动量通量是

\[I>u\]

2 dy 流量相同的理想位流的厚度等于 8-8 1 ，因此它的动量通量是

\[[V~( \delta -81 )\]

两者的动量差，或者说由于粘性使流入边界层的动量通量和位势流相比损失了 z 川 (8-81 ) - I>u 2 dY 已知 U， (8-8 1 )= I: udy. 所以动量通量损失可写为 川 ($\delta$-81) - I> 内=d: U呻 -4:"y =J:川 U， 一山 边界层内损失的动量通量相当于理想流体以速度 U， 流过厚度品的动量通量 p82 U! , 令两者相等可得 4. 边界层近似的推广 (1)曲面边界层

\[82=jf t(I-t)dy\]

上面以平壁绕流为例说明边界层的概念和导出边界层方程，它可以推广到曲壁 绕流的情况，只要壁面的曲率半径远大于边界层的厚度，或者说壁面的曲率半径和绕 流的流向尺度属同一量级.对于曲壁边界层，边界层方程中的流向坐标 z 应是沿曲 壁的曲线坐标，边界层方程中的横向坐标 y 应是垂直于曲壁的法向.边界层方程中

\section*{10.2  不可压缩流体层流边界层的相似性解}\addcontentsline{toc}{section}{10.2 不可压缩流体层流边界层的相似性解}

的 U，(.r) 是绕曲壁元黠流动的壁面速度. (2) 可压缩边界层 大 Re 可压缩绕流也可利用边界层概念导出边界层近似方程，近似的要点是:${\rm 1}$边 界层内压强在垂直壁面方向不变，沿壁面方向的压强等于外部无粘流场的当地壁面 压强 z ${\rm 2}$粘性项只有法向扩散，流向的粘性扩散远小于法向扩散，可以忽略，即在可 压缩流体的纳维 - 斯托克斯方程中忽略)1 iJ 2 U / ${\rm e}$1.J. .2 和 d与/ J .r2 ${\rm •}$ 考虑可压缩流体边界 层内密度将有变化，如边界层内的密度为 $\rho$ (y) ${\rm •}$ 外部位流在物面的密度为户，则根据 质量流量的连续性原理，在进行同样的推导后可得可压缩边界层中的位移厚度为 81 = | 《们 (1 - PTUT )dy (10.15) J ()飞 p. $\upsilon$ ， I 同样推导方法可得可压缩流体边界层的动量损失厚度的公式如下 z 2 = I 归 ' \_p\_$\iota$( 1 一 旦 )dy (10.16) J 0 p ， U ， 飞 u， r J (3) 非定常边界层 非定常绕流也存在边界层，这时边界层的厚度、壁面的摩擦应力等都是时间的函 数.只要非定常流的时间尺度 T- L/ U . . 在定常边界层的流向运动方程中加上局部 加速度项，就可得非定常边界层方程. (4) 薄层近似 普朗特边界层理论的核心部分是绕流的图壁处粘性层非常之薄，将这一分析方 法应用到其他很薄的粘性剪切层，将得到形式相同的方程.因此边界层方程可以推广 到其他薄剪切层的流动.例如棍合层、射流等，所以边界层近似也称为薄层近似.以下 各节将利用薄层近似方程分别研究边界层、 1昆合层、射流等各种流动的特性.

下面讲述边界层方程的解法.虽然边界层方程已经抛物化，但仍然是非线性的对 流扩散方悍， 一 般'情况下很难求出解析解.大多数情况只能求得近似解或数值解.本 节介绍 一 类解析解，称为相似性解.这种方法在流体力学及其他非线性方程的求解中 也是很有用的. 1.相似性解的概念 相似性解是 一 个非常有用的概念，它可以将二维的偏微分方程简化为常微分方 程.我们以平面定常不可压缩层流边界层方程为例来说明这个概念.在 10. 1 节已经 导出有量纲的边界层方程和边界条件:

u ${\rm ••}$ dU. , J$\tilde{u}$ u'\textasciitilde{}-+v 子 = U. $\tau$二 +$\nu$ 二号

\[r1 X oy GX <l y-\]

2旦十字 =0

\begin{gather*}
<IX <l y\\
y = O. u = v = 0
\end{gather*}

y $\rightarrow$$\infty$ ${\rm •}$ u=U.(x) 通常情况下，对于给定的运动粘性系数 $\nu$ 和边界层外边界速度 U. ( 川，上述方程 的解有以下形式 : u/U,= J(x.y) ， J 为 x ， y 的二元函数.在某些边界条件下，上述方 程的解可写为 在 = f( 1J) 00. 17> 式中 q 是 x ， y 的某一特定函数，则式<1 0.17) 称为边界层方程的相似性解，甲称为相 似性变量.在方程具有相似性解的情况下自变量的个数由两个 (x ， y) 减少为一个甲' 因此偏微分方程可简化为常微分方程.这就使求解方便得多.下面首先来讨论什么条 件下存在相似性解.以及如何找到相似变量. 2. 相似性解的条件和构造方法 首先引人线性变换群如下 z

\begin{equation*}.1' = Ao, x. y = AO, y , u = Ao, u , v = Ao, v , U. = Ao;u.\tag{10.18}\end{equation*}

式中 A 称为"变换参数"， $\alpha$) .的哟'酌'比为待定常数，应用<1 0. 18) 变换群，不可压缩 流体边界层原始方程(1 0.9). (1 0.10) 化为 ACa" - 0 ,' i!..旦 + A' o, 0 ,' i!..卫 =0 (J;r J Y A CZ町 Ql) U 己旦 +A 《 az L aA a23 u44=AUez aiV ，年 +ACOj 川 414

\[<lX (1 Y <l X <1 Y\]

若以上两式中 A 的罪指数都相同，则原变量系 (u ， v ， x ， y ， U，) 满足的方程与新的变量 系 (u ， v ， X.y.u ，) 满足的方程是相同的，也就是说方程 00.9) 和 (10. 10) 对于变换群 (1 0.18) 是不变的 .A 的事指数相同需有下式成立 z $\alpha$3 一-$\alpha$)=$\alpha$1-$\alpha$2 2$\alpha$3 -at =的+也一的 =2$\alpha$5 一 $\alpha$)=$\alpha$3 - 2a2 由以上等式可得出

\[a3 - \alpha  ;. \alpha )-2\alpha 2 .\alpha  ,  = - aZ (10. 1 9 )\]

利用式(1 0.19) ，可在变换群(1 0.18) If1 消去变化参数 A. 从而求而变量系 (u.v ， X.y ， U，.) 与变换系(丘 ， v ，.1' 巾 ， U. )之间的关系如下: A = (.1'/:1.') 才 = (y/豆时， 或 y/XfJ = y/x${\rm ß}$ (10.20) 式中 ${\rm ß}$= $\alpha$2 / $\alpha$). 类似地可由(1 0.18) 的其余各式和式 00.19) 得到

u ${\rm U}$ v v U, U, p习 -FTEE'J7-Z7$\cdot $F=EE-p=EE (10.21 ) 显然，式 00.20\} 和式。 0.2 1)所表示的变量组合是线性变换群式(1 0.18) 的不变 量.因此称为"绝对变量"如果边界条件(1 0.1 1)和 00.12\} 在相应的变换之后不增加 新变量，也只和绝对变量有关，则上述绝对变量就是相似变量.于是我们可设 平=玉， j(q)=J苟. g( 守\} =刀 ${\rm •}$ h(.r\} = 主运 00.22\} u

\[.1.'Y .r' -,' .r " .r\]

式中 f.g 只是平的函数.下面将证明:对于有相似性解的流动 .h (x) 必须是常数.将 边界条件(1 0. 11) 、式 00.12\} 用式 00.22\} 所定义的绝对变量来表示.可写作 U, (x\} TJ=O , f=g=O; TJ- $\infty$ ${\rm •}$ f( $\infty$)=?口T 若有相似性解，边界条件的表达式只能和 q 有关，因此，最后一式必须与 r 无关，即必 须满足 U.( .r\} 一?寸..;... = h(x) = const. Z 俨 令 l 一年=m ， 贝IJ

\begin{equation*}U.(x) = cx"\tag{10.23}\end{equation*}

就是说:只有当无粘外流的内边界或边界层的外边界的速度是 r 的事画数时，不可 压缩流体边界层才有相似性解，前面导出的绝对变量成为相似变量.将式 00. 21 \}和 式 00.22) 代入方程( 10. 9) 和 00. 10\}. 就可将原偏微分方程变换成如下的常微分 方程 z mf 一气旦矿'+g' =0 00.24) mf2 一气勾I' +gf' = 阳 2+$\nu$1'" 00.25) 式 00.24 )和式 00.25) 是函数 f.g 对于自变量甲的常微分方程. 将以上相似性解的导出过程再简单地归纳如下. ${\rm 1}$定义线性变换群并将其代入原始方程 z ${\rm 2}$令每一方程对于变化群是不变的，从而得到变换常数之间的相互关系; ${\rm 3}$消去变换参数，得到绝对变量; ${\rm 4}$用线性变换群代人边界条件，看它们是否只与绝对变量有关 z 若是，则绝对变 量就是相似变量;若杏，则该问题无相似性解. 3. 相似性解的流动特征 具有相似性解的边界层流动，其速度剖面具有相似性.以不可压缩流体平面定常 层流边界层流动为例，由式 00.22) 可得 丘=\_!\_.\{2'.\_\textbackslash\{\} 一卫一= \_!\_\textasciitilde{} \{ 2'.\_ \textbackslash\{\} U. CJ \textbackslash\{\}x P 尸 U..1.,p ' l C 岳飞 xPJ

上式表明，任何 z 坐标处的速度剖面可用单一曲线 u/U，= f<-'l )/C , v/u,:cP 1 =g( 泸/ C 来表示，其中.，，= y/ :cP. 就是说:只要将横向坐标缩小 :cfl 倍，则流向速度 u/ 队的横 向分布相同，即速度剖面的形状相同 z 对于垂向速度 v. 除了缩小横向尺度外，速度值 也要缩小 :CP - 1 倍.这种速度分布称作仿射相似的速度剖面. 4. 半无限长乎板定常层流边界层解 布拉休斯 CBlasius) 解 (1)布拉休斯解 应用上述方法，求解半无限长平板上不可压缩流体平面层流边界层内流动的相 似性解.假定平板为半无限长，平板外无粘流场是均匀的.因此边界层的外流是 u.=U . . . . 为了求解方便，先引人平面流场流函数1${\rm þ}$': J 1${\rm þ}$' cJ1${\rm þ}$' u=--. v= 一一-一 rJ y cJ:c 则方程和定解条件可归结为 J 1${\rm þ}$' J2 1${\rm þ}$' J 1${\rm þ}$' cJ2 1${\rm þ}$' cJ J 1${\rm þ}$' -一 ------ --- = $\nu$-二二

\[cJ y J:cJy J:c r1 y< a yj\]

J 1${\rm þ}$' .. \_ J 1${\rm þ}$' :c >O.y=O:-::-一 =0 或 V=O Z7一 =0 d :C uy cJ1${\rm þ}$' \textasciitilde{}. :c> O.y $\rightarrow$$\infty$. -一 = U , = U rJ y ${\rm 1}$用群论方法求出相似变量 首先引入线性变换群 : :c =A.I :C. y=A.2y. 1${\rm þ}$'=A.,,1${\rm þ}$' ${\rm 2}$将线性变换群表达式代人流函数方程，得到 J l( 壁主主一壁些) = A.，， -3川堕 Jy Ji: Jy cJ.T iJ.V2 , cJ yJ

\[- J 1li'\]

i: >O ,y=O: 1${\rm þ}$' =O. -二= 0; r1 y J 1${\rm þ}$' :c >O.y- $\infty$ :A.' .2 -一 = A..U队， rJ y ${\rm 3}$求指数间关系 若方程对于变化群不变，则必须有 2$\alpha$3 - al - 2a2 = $\alpha$:1 一 3$\alpha$2 .向 - a2 = a4 由以上两式可解得 $\alpha$ :1$\alpha$l 一句， $\alpha$1 -$\alpha$1- 2$\alpha$2 这就是说.对于线性变换群的绝对不变量是 1 A = (; r\textasciitilde{} = ($\tilde{f}$ = (;r; = (兰)\textasciitilde{} ${\rm 4}$消去 A. 求变量组合 z 令的 /$\alpha$1= 卢$\cdot $则 $\alpha$ :1 /$\alpha$1 =l-${\rm ß}$.a 'l/$\alpha$1 =1-2${\rm ß}$ 有绝对变量

主 JA-主川 '1[1' '1[1'

\[zl p-FE'\]

因此可设 "1= 乡 $\cdot $ j( "1) = 主 ${\rm 5}$确定相似性变量 将以上绝对变量代入边界条件表达式，显然，当 ${\rm ß}$=1/2 时，边界条件不增加新的 变量(与前面的结果一致) .所以绝对变量即为相似性变量，即 "1= y/.[i. j( $\rho$= 川使边界条件无量…们可以取如下相似性变量 : "1= 在.j( 甲)= 7主=-ft 入原方程和边界条件，可得

\begin{gather*}
,JJI.T U <.\\
2f"'+ j(' = 0
\end{gather*}

可= o. j = 0 , (= 0; y $\rightarrow$ $\infty$ ${\rm •}$ J' = 1 式中上标 "f" 表示对市的导数.所以，半无限长平板上不可压缩流体定常边界层流动 归结为求解上述 三 阶常微分方程的边值问题，该方程没有解析解，只能用数值积分法 得到该边值问题的数值解. 1908 年，布拉体斯首先得到了这个形式的解，故通常称为布拉体斯解.后来 1938 年霍华斯( Howarth) 用数值方法精确地得到了这个解的数值结果.表 10. 1 给出了计 算结果. 襄 10.1 布拉休斯解的结果 '1= y ./fT7百 f( 守) (<甲) = u /【j 的 j气 '1) 0.0 0.0 0.0 0.33206 0.4 0.02656 0.13277 0.33147 0.8 0.10611 0.26471 0.32739 1. 2 0.23795 0.39378 0.31659 1. 6 0.42032 0.51676 0.29667 2.0 0.65003 0.62977 0.26675 2.4 0.92230 O. 72899 0.22809 2.8 1. 23099 0.81152 O. 18401 3.2 1. 56911 0.87609 O. 13913 3.6 1. 92594 0.92333 0.09809

续表 1J =Y ./fJ万百 f( '1) 1'( 甲 )=u/u /'( '1) 4.0 2.30576 0.95552 0.06424 4.4 2.69238 0.97587 0.03897 4.8 3.08534 0.98779 0.02187 5.0 3.28329 0.99155 0.01591 5.2 3.48189 0.99425 0.01134 5.6 3.88031 0.99748 0.00543 6.0 4.27964 0.99898 0.00240 6.4 4.67938 0.99961 0.00098 6.8 5.07928 0.99987 0.00037 7.2 5.47925 0.99996 0.00013 7.6 5.87924 0.99999 0.00004 8.0 6.27930 1. 0000 0.00001 8.4 6.67923 1. 0000 0.0 (2) 不可压缩平板边界层的特性 下面用布拉休斯解(根据表 10. 1 的数据)分析平板边界层内的速度分布和摩擦 阻力等特性. ${\rm 1}$边界层内的速度分布为 圭'-:. = f'( $\psi$ <10.26) 最俨严 =$\div$M('〉 -f( 市 )J (10.27) 根据上两式及表 10.1 中函数 f( 市〉和 f( 沪的数据，可分别画出边界层内两个速度分 量的分布曲线于图 10.3. 边界层内速度横向分量和流向分量之比为 v 1 1 $\tilde{U}$,., x z\textasciitilde{} 万TZ可 .j言 . l(er = 丁 当 y 不断增大时，横的速度逐渐增加利用表 10.1 的数据可 f\textasciitilde{} ， 1J=8. 4 时，旦俨= O. 邸，精确的极限值为 y $\rightarrow$$\infty$ :V = 0.865 \_\_!主= JRe.

立 JE手 图10.3 平板边界层内速度分布 实线 : U; 虚线 : V 这就是说，在边界层外缘，流体质点有一个微小的垂直主流方向的速度，它对外部势 流场有一个排挤作用，由于 u 相对于 u 是高阶的小量，在边界层内，速度分量 u 常常 被认为是次要的.如果要更精确计算边界层对外部无粘流的影响，则可在更高一级的 近似中计及边界层的排挤作用. ${\rm 2}$边界层的厚度 由图 10.3 可看出，随着 y 的不断增加 .u 渐近地过渡为 U \textasciitilde{} ${\rm •}$ 若规定 u=O. 99U .\textasciitilde{} 处的 y 坐标值为边界层名义厚度 8. 那么，由表 10. 1 用内插法可求出 币 1. ' l1 < =o.\textasciitilde{}!J = 8${\rm Æ}$ = 4 叫 于是，平板边界层的名义厚度为 8 = 5. /旦 =\textasciitilde{}$\iota$ ."$\upsilon$ 骂 、IRe ， ${\rm 3}$边界层内的剪应力和表面摩擦阻力 边界层内的剪应力 z T.r.v -省 =$\mu$JZf( 币) 平板表面的摩擦应力: Tu.. = $\mu$( 去)户。 =$\mu$ 存('(0) 由表 10. 1 可得 \{'(O) =0.332. 因此壁面的摩擦应力 z 日 =o m jdL 0332EL ." X JRe. 壁面摩擦应力常用无量纲摩擦阻力系数表示.

\[<10.28)\]

定义 10.3 局部摩擦阻力系数 CJ 等于局部壁面剪应力除以边界层外的动压

\[(pU; /2) :\]

C r 一 r主

\[pU~/ 2\]

将剪应力计算结果代入上式，得平板表面局部摩阻系数: c, = \_:'U' .\_ = 0.664 一 一 广 声寸2 ../Re: 由式 00.29) 可以看到，平板表面的摩擦系数都与JX成反比. ${\rm 4}$平板上的总摩擦阻力 00.29) 由壁面剪应力公式可见，当 z $\rightarrow$ 0 时 'LU' $\rightarrow$ $\infty$ ，这显然是不合理的，其原因是:当 r$\rightarrow$0 时，边界层简化前提不成立.严格地说，在边界层前缘需要用纳维 - 斯托克斯方 程直接求解.如果平板很长，边界层前缘的剪应力对整个平板的摩阻贡献很小，因此 仍可应用边界层解的结果来计算平板上的总阻力.长度为 L 的平饭其单位宽度上所 受的总摩擦阻力为 D = j; 凡 d .r = O. 664 ;;P盯 总摩擦阻力也常用无量纲总摩擦阻力系数表示. 定义 10.4 总摩擦阻力系数等于总摩擦阻力除以摩擦面积和边界层外的动 压，即 C D 一 广严$\tilde{L}$ / 2 将不可压缩流体绕长度为 L 的平板定常边界层的总阻力结果代人，得

\[C" = 1. 328 -\]

0 万言 其中 Re l， =U.. L / $\nu$. (3) 速度剖面相似性的验证和布拉休斯解的局限性 00.30) 实际平板绕流总是有限长度的，采用无限长平板的相似性解来近似，其准确性需 要实验验证 .1942 年尼古拉茨发表了零攻角平板边界层内速度分布的测量结果，他 在风洞实验中得出了与平板前缘不同距离上的速度剖面，并按相似性变量 (U / U . , y ..;u二万万整理数据，其比较结果如图 10.4 所示.实验结果很好地证实了速度剖面 的相似性，也验证了布拉休斯理论解的正确性. 最后，讨论布拉休斯解的应用范围及修正.从这个问题的建 立 和求解过程来看， 布拉休斯解受到 三 个假设的限制: ${\rm 1}$作边界层简化时，假定 I J 2

\[u / a .r\]

2

\[1 << la\]

2 u/ay 2 1. 在平板前缘附近，这个假设显 然不成立，因此布拉休斯解不适用于平桓前缘，事实上，前面已经发现布拉休斯解在 前缘有奇性. ${\rm 2}$写边界层外缘边界条件时，假设边界层外的流动不受边界层存在的影响，事

1.2 1.0 I::J 0.8 0.6l 9恨;正 + 0.85 口 0.86

\[0.4| L3FV\]

。 0.93 A 0.82 x 0.93 0.2$\iota$-JJ早二 口 1.08 ${\rm •}$ 1.1 1 A 1.24 2 3 4 5 6 7 8 FYJE 图 10.4 布拉休斯理论解与尼古拉茨实验结果的比较 实上边界层外边界有法向速度，虽然是高阶小量. ${\rm 3}$假设边界内的流动不受平板后缘的影响，或者说，平板是半无限的.在有限氏 平板的后缘附近，这个假设显然不能成立. 为了考察上述 三 个条件对布拉休斯解的影响，郭永怀从准确的运动方程和边界条 件出发，求得了绕有限平板流动的 二 阶近似解.他给出了总摩擦阻力系数的修正公式 z

\[c" = 1. 328\iota 4.12\]

广亏京-， ReL 式中下标 L 表示 Re 数的特征氏度是平板长度 L. 实验结果表明，布拉休斯解适用于 ReL > 100. 在 ReL::::;;;;100 范围内应采用郭永怀 二 阶近似解. 5. 平面层流射流 没有固壁的粘性剪切层称为自由剪切层，无界射流、绕钝体流动的尾流和混合层 均属于自由剪切层.下面研究平面自由剪切流的薄层流动，首先考察平面射流. 设流体从无限长狭缝向充满同种流体的空间喷 出，形成平面射流，如图 10.5 所示.该流动由喷流驱 动，全流场的压强是常数.由于粘性作用.周围流体 被喷射流体卷吸，从而使射流的宽度向下游逐渐扩 大.实验证明，当雷诺数足够大时，喷射流体和周围 流体的混合区域是 一 个薄自由剪切层.若狭缝宽度 无限小，可将坐标原点取在狭缝出口平面，设定主流 方向为 z 轴 .y 轴垂直于主流方向，流动对 y 轴对称. 图10.5 平面射流示意图 x

(1)基本方程 作为薄层近似，这一问题的流动控制方程和零压强梯度的层流边界层方程相同: (2) 边界条件 ${\rm 1}$对称条件 ${\rm 2}$渐近条件 当+些 =0 o.r "y au , dU i) 2 U

\begin{gather*}
U -=-- ;- v -=-- =\nu --\\
\theta  .r ' ~ () y \tilde{dy}2
\end{gather*}

y=O. V=O , 旦旦 =0ay y$\rightarrow$土 $\infty$ ， U=O 很显然 U= lI =O 满足上述方程和边界条件，但不是所要求的解.事实上，流体由 狭缝喷出，还需要给出射流的人口条件.本问题中，不可能给定人口的速度，因为我们 假定入口射流是一个点.但是可以给定人口的射流动量通量 ] ， 因为全流场压强梯度 等于零，无界射流流场在没有其他任何外力作用的情况下.通过射流任意横截面的动 量通量 j 应等于人口的动量通量.于是可以补充积分条件:

\[] = d ' - \tilde{u}2 dy = \infty  00.3 1)\]

补充这个条件后，该问题就有唯一解了.与上题同样方法，引人平面流场流函数 $\psi$: cJ'IJ! ${\rm e}$J'IJ! U= 一一-， v=---

\[<7 y d.l'\]

则平面层流射流的方程和定解条件可归结为 a'lJ! d 2 'IJ! d 'IJ! () 2 'IJ! J 3 'IJ! -一一----- -一- -$\nu$--ay ${\rm e}$J.ray a.r ${\rm e}$1 y ' ay 3 ${\rm e}$1'1J! \_ a2 'IJ! .r >O ,)I =O: \_- =0 ，一一= 0; v $\rightarrow$士 $\infty$ :一一 =0 V ${\rm •}$ (J.1' . . ay 2 二 V - ${\rm •}$ ay

\[T (.,.. , a'lJ!   2\]

.r> 0: -=- = I 1-:::-1 dy = const. P J -" 、 rJ y I (3) 相似性解 (1 0. 32a) (1 0. 32b) 00. 32c) 注意到这一流动问题的边界条件和边界层流动类似，在流向没有特征长度，因此 该方程和边界条件可能具有相似性解.仍用前述群论方法求出相似变量和相似性解. 引人线性变换群 z .r = A", .i, y = A"2 Y ， $\psi$= A"l 苦 (10.33) 和导出层流平板边界层相似性解的方法相同，将式(1 0.33) 代入式 00.32) ，消去变量 A ， 可求出指数间的关系为 2 1 $\alpha$ z= 言剧， $\alpha$3 = TQI

\noindent\textbf{10.2.} 不可压缩流体层;在边界层的相似性解

这就是说，对于线性变换群(1 0.33) 的绝对不变量是

\[y - y\]

p$\pi$ - :i 2 万'

\[X 1/ 3\]

根据上面分析，可选定如下变量组合为相似变量 z 'I${\rm þ}$' .i 113 市=打在万 ${\rm •}$ f( 守) = -;;-:-I毛古0"'''. x.'. , 2qv 一;::: .1" 1 一 (10.34 ) 式中 q 是速度量纲，引入常数:.rv和 q 使币和 f 成为无量纲量，也可使方程在形式上更 简单.由式(1 0.34) 可得

\['I\theta ' = 2qv 1/ 2 X 1/ 3 f\]

应用链式求导法则主='!!1旦旦 =2251 ，可得 iJx i1 xd 守 ${\rm •}$ ay i1 y d 守 rJf = 一三豆兰丑， EI= $\rightarrow$手$\tau$7$\tau$(($\eta$ 〉i1x 3 x ay 3", ' / <X< / ;!J '" ;3=( 萨和 )zf'( 训， 去 = (3 "， 1/ 年言 f 阳市〉 射流中的速度可由上式导出 z u=21雪 =?rnal/Zzl/3EI= 主，， 2 五组2 3y"哇" $\tilde{i}$1 y 3 哩 X 1 / 3 v= 一古=一号 4[州一 2 rJ\{' (7J) ] r1 X " J.. 将以上各式代人式(1 0. 32a)- 式(1 0. 32c). 可得

\[j"" + 2( 1'2 + f f") = 0 (10. 35a)\]

平 =0. f=I'=o. 市$\rightarrow$士$\infty$ ${\rm •}$ l' = 0 00, 35b) 现在射流的相似性解化为以上三阶常微分方程的边值问题，式中待定常数 q 可由积 分形式的条件 jl : 向=岭来确定 下面求解方程 (10.35 剖，将它帜分一次，并应用甲 =0 时的两个边界条件得到 z

\[1'+2ff' = 0\]

再积分一次得

\[l' + P = const.\]

考虑解中有待定常数 q. 可取积分常数 const. =1. 这相当于认为 1' (0)=1. 所以上式 写为 1'+ 尸= 1 再积分上式，并注意 f(O) =0. 可得到 7J = r\textasciitilde{} 1 \textasciitilde{}乌 =iln 过工=盯thfJ" 一 F 2 ---1- f 00. 36a)

或

\[1 - ",-2,\]

f= 问 =t仨7 下面用射流动量通量的积分条件来确定常数 q. 代人速度 表 达式: u= 毛毛f'( 亨) =号$\rho$I 凡 (1一 th 2 1J).) x . - .) 将上式代人式(1 0.3 1)的动量通 量 守恒条件得 f=j: 向 =fdq3j ti (l-M 州 =pq3 由此可得 q = 0.8255(右) 1/ 3 给定 J 以后 .q 就完全确定了，平面射流问题就全部解出了. (4) 结果分析 将式(1 0.36) 和 q 值代人速度表达式便可得到射流的速度剖面 z u=o 叫声lLthv v = 0.5 叫乡 ) 1/:' [刷一 th 2 1J)一问] 式中 00. 36b) (1 0. 37a) (10. 37b) 1J = 0.2752 (兰) 1 / 3 二名言 (10.38) 显然式(1 0.37) 给出的速度分布具有相似性.令守 =O(y=O) ${\rm •}$ 可得到对称轴上的速度 分量 z ul 户。 = Umax = 0.4 叫乒 i1 /3 .u| 户。 =0 飞 P- J.IT I 通常，可将 z 方向的速度分量写成

\[~=1-th2 1l\]

U ma粟， 上式给出的元量纲速度剖面见图 10.6. 射流外边缘的速度分量为 ul 户 = O. v 1..\_'" =一 0.550( 句 l i 3 \textbackslash\{\}严. - I v( $\infty$ )<0 说明无穷远处有流体进入射流，称为射 流的卷吸作用，它是通过流体的粘性将原静止流体 卷入射流中来，通过射流横截面的单位宽度上体积 流量可通过积分求出 z 0.8 \textasciitilde{} 0.6 3 ::: 0.4 0.2 11 L 过

\[-4 -3 -2 -1 0 1 2 3 4\]

可 图 10.6 射流速度的相似性分布

r+'" I ",, \textasciicircum{}IJ \_\_\textbackslash\{\}I ' l Q = I ud .y = 3. 30 (一$\nu$:r) (10.39) J ，. 飞 p , 由上式可看出，体积流量 Q 是沿 z 方向逐渐增加的，这是射流不断卷吸周围流体的 结果.同时可看出 Q 与射流的动量通量 J 的立方根成正比. 由阿孔喷出的轴对称射流，也有相似性解.其求解方法和平面射流问题类似，可 参见粘性流体力学的专门书籍，如希利希廷 (Schlichting) 的经典著作《边界层理论)). 6. 两平行流动间的层流剪切层 在许多化 t 和燃烧器 11 1 常常有两股平行流动的汇合，在自然界巾也常有两股平 行气流汇合的现象.当两股速度不同的流动互相接触时，由于粘性作用.在它们之间 会形成 一 个与边界层类似的剪切层.当这两股流动的特征雷诺数足够大时，它们间剪 切层很薄.且横向速度比纵向速度要小得多.如图 10.7 所示，该问题也属薄层流动， 可以用边界层方程求解. (1)基本方程 将实际混合层流动简化为图 10. 7 所示的力学模型，设有两股不可压缩平行流动 以不同速度平行汇合.流体密度为 p. 流体的运动粘性系数为 $\nu$ ，上半平面的流体速度 为 U 1 ${\rm •}$ 下半平面的流体速度为 U2 (U 1 > U 2 ). 剪切层白。点起逐渐发展.这两股流动 间国剪切应力作用.相互间有动量交换.但无质量输运.因而两股流动的分界面是零 流线 '[1" =0. 零流线将流动分为上下两 l 天 .11豆和 H 区，如图10.8 所示.以 0 点为原 点，取.\{"轴平行于来流方向 .y 轴垂直于来流方向. l每I 10.7 两股流动的汇合 一 平面剪切后 图10.8 平丽的切层的流向演化 和平面射流相同，零 Jk 强梯度的薄层方程适用于混合层.由于月二力梯度等于零， 可写出 I 、 H 两区都适用的流函数方程(参见图10.8) : (2 )边界条件 ${\rm 1}$渐近条件 ${\rm i}$J'[I" a2 '[1" .J'[I" .J \textasciitilde{} '[1" CJ" '[1"

\[.J Y ,1.1"<] Y ,]." a y' a y'\]

c7'[1" \_. ,1'[1" y- ' 且 ，一一 = U1 ; y $\rightarrow$一$\infty$.一一 = U2

\[cJ y - dy\]

00.40) (10.41a)

${\rm 2}$交界面上连接条件:设两股流动的交界面方程为 :y=y'( .r)， 界面上有 y = y' : U .I. = U 句 ，叫= v \_ ; $\rho$t=$\rho$ ， T.cy i = T.cy- 00.41b) 这里下标"+"和"一"分别表示交界面两侧 I 区和 H 区的参数.和射流问题类似，平面 混合层流动问题也没有特征长度，因此方程 00.40) 对于边界条件(1 0.4 1)存在相似 性解，具体推导过程如下.对于 I 区可作如下相似变换 z '71 = y..../哇 ， 1Jr1 = j ,v,.r 1, ( 币，) (10$\omega$ V !CC 由此可得

\[\cdots 11. ( '71 ) , v , = ~ ~ ['711. ( '71 ) - 1, ( '7' ) ]\]

咱 1=$\mu$ 穹 =PZf;( 甲) 利用相似变换可得到常微分方程 z 2 广;+ 11f\textasciitilde{} = 0 对于 E 区可作同样相似变换 z 加 =y 在，矶=而叫 由此得 均=川)， uz=$\div$ 汪叫2 )叫 )J $\tau$呐 =$\mu$ 穹 =p 冉f州 且得到常微分方程 00. 42b) (10. 43a) 2 产;+ 121飞= 0 (10.43b) 这就是说元 ， /2 分别为'7' ，如的函数，且满足的方程相同.通过上述变换，边界条 件 00. 41a) 可写作 z \{川$\infty$ z 川=: 加$\rightarrow$一$\infty$ z 儿〈加 \}=t=$\lambda$ (1 0.4$\omega$ 此外在交界面上''71 = 币，和 =rJi( 对于同种流体的汇合 :rJi l '7 t =1); 由联结条 件(1 0.40b) 可得 (川) =斤〈而) = 0 11 句t)=/\textasciitilde{}( 守2') ， f\textasciitilde{}( '7 t)= f'z("I2') 为了消去联结条件中的 '7 t 和币，重新定义自变量，令 " = '71 - '7t , '2 =和 -"12' 00. 44b)

\section*{10.3  卡门动量狈分关系式}\addcontentsline{toc}{section}{10.3 卡门动量狈分关系式}

用，.和 '2 作为自变量，混合层流动相似性解的方程、边界条件和联结条件可写作 z 2 户;+f.1才= 0, 2 月 + f2 ${\rm I}$z = 0 00. 45a) VI( $\leftarrow$ $\mu$0) = 0 ，只($\leftarrow$川 1\textasciitilde{}(o)= ${\rm I}$z (O) ， ${\rm I}$I(+ $\infty$) = 1. 只(一$\infty$) =,\textbackslash\{\} 00. 45b) 其中 '\textbackslash\{\}=UzjU. 注o. (1 0.45) 各式是关于自变量，.和岛的常微分方程组，这组方程 没有初等画数表示的解析解，可用数值计算的方法以及职分关系式的近似方法求解， 流向速度分布如图 10.9 所示，它可近似为 Un = $\tilde{U}$2 -U. )th 守 +U2 -U. ., 2 2 E 's' 4 3 2

\begin{gather*}
-2\\
-3
\end{gather*}

4 f o 0.25 0.5 0.75 I 1.25 1.5 吨 图 10.9 平面混合层的相似性速度分布

虽然边界层方程是纳维-斯托克斯方程的近似方程，但它仍是非线性偏辙分方程， 对于不可压缩牛顿流体来说，也只有少数几种流动可以找到相似性解，如 z 绕平板或绕 模形物流动等.一般情况需要用级数方法或直接数值方法求解，但工作量是相当大的. 在实际工程问题中，需要有简便的方法能够较快地估算物体表面摩擦.边界层理论的卡 门动量积分关系式就是一 种积分近似法，它不要求边界层内每一点都满足边界层方程， 只要在积分意义上满足边界层方程，这种近似方法在工程中得到了广泛的应用. 1.卡门动量积分关系式 下面推导动量积分方程，不可压缩流体平面定常边界层方程为 些+争 =0 OI ny u. au \textasciitilde{} $\tilde{dU}$\_ . a2 u u 王三 +v ::，- = u, -;--二 +$\nu$7号(/:r n y a:r d y-

将第一式乘 u 和第二式相加得 dU. , rJ 2 U 一一+一\textasciitilde{} = u. \textasciitilde{}:--二十 v 一$\tau$dX 'dy - r dx ' - ${\rm e}$J y${\rm ‘}$ 再用 U， (X) 乘连续方程可得 以上两式相减后对 y 积分得 ${\rm e}$J uU, , rJU,v dU, 一十一-..--

\[Jx Jy - dx\]

j $\iota$ 立问，一山y+ 问，一 u) ] I \textasciitilde{}

\[o\theta r\]

dU, r""r " r7 ul +寸一 I (U. - u)dy = 一 $\nu$$\tau$| aXJ u OYI " 边界层流动的边界条件有

\[y=O:u=O.v=O; y=8: u=U,. Ju/ rJ y=O\]

因而以上积分式中 [v(u , - u)J I\textasciitilde{} ' = o. 一 $\nu$ 争 l ' = v\textasciitilde{}旦| "Ylo OYI 将以上各式代人积分式可得 d r. ,... ${\rm •}$ y ${\rm •}$ ${\rm •}$ -,. , dU \_ r-'- ${\rm •••}$ ${\rm •}$ ${\rm •}$ J u I \_$\tilde{I}$ [u(U.-u)Jdy+u\textasciitilde{}二 I (U, - u)dy = $\nu$$\tau$| aXJn aXJo OYI...o 而 $\mu$ dU/dyly$\tilde{o}$= 凡，故上式可写为 d r-'- r "r ,--", dU, r \_$\tilde{I}$ [u(U. 一 u)Jdy 十 U\textasciitilde{}' I (U, - u)dy = lu' aXJo aXJII p OO.46a) 根据位移厚度和动量损失厚度的定义式(式(1 0.13) 、式( 10. 14 门，式 (10.46a) 又 可简化为 ? \textasciitilde{} ${\rm •}$ , ${\rm ••}$ dU. \_ T 丁一 (U\textasciitilde{}82 ) + U, \textasciitilde{}:--二 8 1 二 2 ax ax p 引进形状因子 H=81 / 仇，上式可写为 d82 , 82 , n , .n dU, Tw 一一 + :'; (2 + H) \textasciitilde{}一= :':. (10. 46b) dx ' U,' - , .., dx pU\textasciitilde{} 式 (10.46b) 积分关系式首先由卡门导出的，所以称为卡门动量积分关系式或动量帜 分方程. 动量积分关系式含有 H.8 和 T u 共 3 个未知量(或者仇 .82 和 Tu') ${\rm •}$ 显然是不封闭 的但是，所有这些未知量都和边界层的速度分布有关，例如 : Tu 由学 (0) 确定 .8 1 .82ay 由式<1 0.13 )和式 00. 14) 算出.如果给定边界层的速度分布 u/U \textasciitilde{} = f( 旷的， 方程 (10.46b) 只是一个变量 $\delta$ 的常微分方程，就很容易求解.所以动量积分法的近似 精度，依赖于假定速度分布 u/U. = f( 旷的和实际速度分布之间的符合程度.通常可 以采用多项式近似的速度分布 z

N f( 币) = L:叫.. TJ = 去 式中的系数由边界层方程的边界条件确定.例如:在壁面上，流体满足粘附条件，代 入动 量 方程后，在壁面速度分布必须满足以下两个条件: y = O:U = O. 引 I = \_!\_学=一旦华ay I Y $\tilde{o}$ $\mu$ ax $\nu$ a x 在边界层外缘.速度分布必须满足渐近条件.它们可表示为 (1 0. 47a) aU (J 2 U a.U

\[y = 8:u = U , (x). .~- =::;-\tau = \cdots  ~\tau = \cdots = 0 (10. 47b) dy dy - dy\]

一般来说.多项式的项数取得越多.速度分布的近似程度越好，实际上最重要的是壁 面条件.和外缘渐近条件中的低阶导数项以上条件中生 I =\_!\_业是壁面上的边 (J y 2 I 卢 $\mu$dx 界 层 动 量 方程，它使所设的速度分布和压力梯度联系起来，尤其是它可以根据顺压梯 度 (dp / d x <O) 或者逆压梯度 (dp / dx > O) 确定边界层速度分布在壁面处的特征.至于 边界层外缘的渐近条件(1 0. 46b). 可根据近似的精度要求，选定近似的阶数.下面用 边界 层 积分关系式来求解不可压缩流体的平板边界层流动. 2. 零压力梯度乎板层流边界层问题的近似解 前面通过求解边界层偏微分方程得到了平板边界层流动相似性解的数值结 果 一一 布拉休斯解.下面用动量积分关系式求该问题近似解.由于平板外无帖势流是 均匀流，所以有 dU， / dx=O. 也就 是 dp / d.r=O. 于是积分方程 (10.46b) 可简化为 d8z Tu. dx pU\textasciitilde{} 通常按以下步骤求解边界层动 量 职分方程. (1)确定速度分布 设速度分布为 [pf =/( 沪， rf 选取函数 J. 使它们尽可能和真实速度剖面相吻合，令 I( 币) = L:a. r;" ,, = 0 00.48) 其中 N + l 个系数 a. 确定如下:令多项式满足式<1 0. 47a) 和式(1 0. 47b) 中最主要的 (N 十1)个边界条件，得到 N + l 个代数方程，由此确定岛，以 三 次多项式为例: J< TJ) = 叫 I + hrJz + cTJ + d 该速度分布必须满足边界条件: u l 叫 = U . . ${\rm •}$ 学 I = O. u 1 山= O. I1 YI , ${\rm ð}$ rJ2 U I ay2 I 户。 u

其中前两个方程是边界层外缘的渐近条件 z 第 三 个方程是壁面元滑移条件 z 最后 一 个 是零压强梯度条件.它们可写作: f(O) =0 ,/'(0) =o ,fO) = 1.1 (1) =0 ，由此可得 a= 一$\div$， b=0 ， c=? ， d=o 于是有 3 1 理 U二 = f( 1J) = 言市一言矿 (10.49) 选定了 f( 沪的逼近函数，并不意味着速度剖面完全确定了.因为在 q 中还包含边界 层名义厚度 8 ，它是 z 的未知函数.也就是说式(1 0.49) 表示 一 个带有参数们川的速 度剖面族，为了完全确定速度剖面，需要由动量积分方程确定们 x). (2) 8( 川和t"". 的确定 首先将速度剖面函数代人\&公式， 82 = I: 世 (1 一 rl..， )dy = 8I>0- f) 句 =8I(\textasciitilde{} 1J- 拉 )(1 一书+扫 )d 1J =朵8(x) 然后根据壁面摩擦应力表达式可求得 7=$\nu$( 去) y-o =$\nu$ 号川=号E 其中 ${\rm ß}$=/( [ \textasciitilde{}一打] I ， $\tilde{o}$ = \textasciitilde{} (3) 最后将 82 代人式 00.48) 可得

\[39 d8 3\nu \]

280 dx 28U\textasciitilde{} 职分上式，并注意到 x=0 ， 8=0 ， 可得

\begin{gather*}
8=464JE\\
t"w = o. 323.oU ~,  ./n
\end{gather*}

lJ = 0.323 PU\textasciitilde{} 一

\['v U = x ./Re:\]

所得结果与布拉休斯精确解形式一样，只是系数略有不同，布拉体斯解的壁面剪应力 公式中系数为 0.332 ，动量积分法的系数等于 0.323 ，两者相差甚微.

\section*{10.4  边界层内的流动与分离}\addcontentsline{toc}{section}{10.4 边界层内的流动与分离}

1.边界层内的流动及分离现象 以上介绍了普朗特边界层理论的思想，导出了边界层方程并介绍了一些解法.现

在我们进一步考察边界层内的流动，绕流物体的表面上速度等于零，因此物面附近边 界层动量方程为 (旦旦 L= 工业(J y2Jo fl dJ 式 q:J (J2 U/ ${\rm e}$1 y2) 。是边界层速度分布在物面的二阶导数，几何上说，它是速度型线在 壁面处的曲率.它和 d$\rho$/dJ 成正比.根据边界层流向压强梯度 dp / 巾，也就是边界层 外位势流动的压强梯度，将边界层内的流动分成三种情况，如图10.10 所示. (1)流动方向原强减小，即 dp/dJ<O 的情况，这种情况称为顺压梯度区，此时 (J2 U/ (J y2) " < O. 这种速度剖面的形状是外 凸的(图 10. 10 中 1.2) ; (2) 当压强达到极限值 H才，即 d$\rho$/dJ= O. 称为零压梯度，此时 ((J 2U/ (J y2) () =0. 边界 层速度剖面 u(y) 在壁面上形成 一 个拐点 (图10. 10 中 3); 图10.10 边界层内的流动示意图 (3) 流动方向上压强升高，即 d$\rho$/dJ>O 的情况，这种情况称为逆压梯度区，此时 ((J 2 U / (J y2) 0> 0. 壁面附近边界层的速度剖面 u(y) 是内凹的(图 10. 10 巾的. 下面进一步分析这三种流动情况下贴近壁面的流动.流体质点在运动过程中受 到两个力的作用，一是受到粘性力作用，它使流体质点减速，二是受到压强梯度的作 用.在第一种情况下 (dp/dJ<O) :沿流动方向压强减少时，压强作为驱动力将抵消一 部分粘性摩阻的作用，使流动不至于减速很大，壁面附近的质点仍可以向下游运动; 在第三种情况下 (dp/dJ>O) .流体质点受逆斥梯度作用，这时粘性力和压强梯度都 使流体质点减速，甚至可使壁面附近流体质点倒流(图10. 10 中 5) . 因此，边界层内的速度分布和流向的压强梯度有密切关系.当 dp / dJ<O 的时 候，整个边界层内的流体质点沿正 r 方向运动，边界层内的速度分布为 du/dy>O; 壁 面附近 (J 2 U/ r7 y2)O<0 ， 边界层外缘附近始终有 (r7 2 U /cJ y2 )0<0. 因此边界层内的速度 型线是外凸的.当 d$\rho$/dJ=O 时，边界层速度型线在壁面上 (${\rm e}$\} 2 U / (J y2 ) 0 = 0 .就是说边 界层速度型线在壁面上有拐点，但是在壁面以外仍有 (J 2 U/Jy2) O <0 ， 因此速度型线 仍保持凸的，并且边界层内的速度分布 du/dy>O. 当边界层内开始有逆压梯度时 (dp/dJ>O) 壁面上 (J 2 U/Jy2)O>0 ， 但在边界层外缘附近始终有 ((J 2U/Jy2 )<0 ， 因此 (]2U 在边界层内部出现拐点，就是说边界层内必有某处-一 (y' ) =0. 在逆压梯度的起始r7y2 阶段，壁面附近的流体质点还保持沿正 r 方向的运动，仍有 (du/dy).y$\tilde{o}$ >0. 然而 $\omega$ 2U/ (J y勺。 >0 使壁面附近速度剖面的形状出现内凹.如果沿流动方向继续保持逆压 梯度，压强和壁面摩擦都使质点减速，那么就有可能在壁面上某一点处用 5 表示，达到 (du/dy)y $\tilde{o}$ =0. 在 s 点以后，近壁流体质点在逆压的作用下向后倒退形成逆流，即

(du/dy)，. $\tilde{o}$< O. 这种现象称为流动从哇面上分离 ， S 点称为分离点 ..S 点后称分离区. 以上就是边界层分离的流动过程，普朗特将 Cdu / dy)v $\upsilon$ = 0 作为 二 维定常绕流 边界日流动分离的判据，也称为普朗特分离判据.分离仄中倒流H:f${\rm i}$:形成涡旋.分离 区产作:涡旋的同 H才.边界民迅速增厚- .从以 1: 分析可以看出，流动分离是逆压棉度和 粘性摩阻综合作用的结果，通常在 JII 页 ml汉是不会发唯分离的.只有在逆 J$\tilde{I}$灭内才有口I 能发什二分离现象.普朗特边界层方程原则 1- 只能适用非分离 l 泛.在分成\textasciitilde{} ${\rm þ}$( 内边界层厚 度大大增加.边界层内速度分世之间的最级关系发吁:了变化.不 i'l: 符合边界层理论的 基本假设.分离区常常伴随 ll\textasciitilde{}1 流并产生不同于理想流动的吓.强分布.在 H 闭物型的定 常绕流'1 1 .分离流动的 n\textasciitilde{} 强合力产生流动阻力.这部分阻力往往大于峨面剪应力的合 力.前者称为 ffi 差阻力.盯?f称作摩擦阻)]. 边界层分离 l泛的流均速度分布.明显地有别于非分离 ${\rm þ}$( 的分布.1:14 此边界层的特 征挝 :()I / $\delta$ ， ()I / () \textasciitilde{} 也有明报变化.可以利川 U ， .= cx 例 的边界层相似 n 解，来号'察不同 压强梯度的边界层流动 .111 = 0 是零压强树H支边界应流动 ;m > O 是 JiI'${\rm I}$iJk 梯度流动 z 111 < 0 是逆压梯度流动.在 m < O 的逆 FI\textasciitilde{} 梯度边界以流动 rll 就可以:1-I.f1l.\textbackslash\{\}分离现象，通 过整理已有的计算结果.得到在分离处，边界民的，\textasciitilde{}\textasciitilde{}参数 H = $\delta$1 /()2 = :\textasciitilde{}. 5\textasciitilde{}4. O. 在边 界启动量积分公式中.常用这一参数来判断流动是百分离. 分离流动不符合边界民理论，研究分离流动应从完整的纳维 - 斯托克斯方挥 出发. 2. 绕流阻力及边界层控制 从前面粘性流体定常绕流现象的分析'1 1 .口 I 以看到绕流物体仁，总的阻力由两部 分细成:摩擦阻力和阳是阻力.摩擦阻力主坚巾和li 性作用产什: .对大 tE 诺数层流流动. 摩擦阻力约与流动速度的 :3 / 2 次 jj 成正比. 1M J I\textasciitilde{} 差阻力是指物体 ${\rm i}$j${\rm I}$j J，日-的 m \textasciitilde{}.虽分布的 合力，阳差阻力与分离阪大小直接相关.骂分离|泛很大 H才. H\textasciitilde{} 差阻力是流动阻力中的 主要成分，因此控制边界民分离是减阻的途径之 一 . 在很多工程问题'1 1 .控制边界层分离十分最\textasciitilde{}:.如果机翼1.: ;去而分离!反过大.将 造成失速.控制边界层分离的方法很多.但不外乎附大类. 一 类是改变物形，控制物面 上的限强梯度，从而尽址缩小分离民;例如采 JH 细氏的流线型物面.另，类是考虑流 动的内部因素.增加边界坛内流体微闭的动址以加强抗逆压梯度的能力. $\tilde{Il}$:在壁面 上吹吸流体.延缓分离.减小分离民.达到减少 Jk 差阻力的效果. I-H T 流动的分离点和l 来流的状态有关，例如1 .机翼上的分离点随来流的攻角而改变.[;1;1 此.问:主点吹气或吸 气的控制方法往往不能适山变攻角的情况.近年来.利用做到传感器 ${\rm i}$WJ 址绕流物面的 流动特性(例如压强或 JI\textasciitilde{} 1.虽梯度) .根据测得的信息，在物面必要的位茸实行流动控 制，这种带有反馈信息的控制方法称作主动控制. 边界层控制是很有意义的实际问题，它是粘性流体力学中的专门课题.

\section*{10.5  可压缩流体定常层流边界层的主要特性}\addcontentsline{toc}{section}{10.5 可压缩流体定常层流边界层的主要特性}

不可压缩流体边界层流动适用于液体和低速气体的绕流问题.高速飞行器或流 体机械内部的气体流动往往速度很高，尤其是高马赫数流动 (Ma>l). 绕流物面的温 度和热流量很大，边界层内气体的压缩性不可忽略，必须考虑气体的密度、温度和物 性的变化.关于高速气体流动的边界层问题是粘性流体力学的主要内容之一，本书概 要地说明它的主要特性. 1.可压缩边界层流动的基本方程 现在来考察常比热容完全气体的平面定常流动，即 r1 \_ a 丰 -..一…$\leftarrow$一… rJ t rJ z 气体状态方程为

\[\rho =\rho RT. R=cp-cv. e=cvT. h=c\theta T\]

气体的分子粘性和分子扩散系数之比称为普朗特数 Pr: Pr = 庄旦= const 为了分析简化起见.假定分子粘性和分子扩散系数都是常数，不随温度变化.在以上 假定下牛顿型气体的二维定常流动服从以下方程组. 连续方程 但旦+组卫 =0 cJ.r . a y 动量方程 u. i1 u CJ b . a , r1 u 、 ${\rm e}$J r ,au , iJ v\textbackslash\{\}l puz+pzy;;=-jJ.+ 石 l f-! J三)+巧 l l-t l 巧+石 )J $\rho$u 去 +P11 艺=一劳 + ad:r[ $\mu$( 去+去) J+ cJay ( $\mu$ 去) 能量方程 pu 去 +PZ=u 去 +v 艺 + ';$\tilde{l}$. ( $\lambda$ 主)+ ;y( $\lambda$ 艺)+ s6 骨是耗散函数: s6 =$\mu$[2 (去 f +2( 号 f 十(去 +ZY] 与不可吓.缩流体边界层方程导出的方法相同，认为壁面法向的速度分量和长度 尺度与流向相比小一个世级，基于这一估计，可压缩流体边界层流动的连续方程和动

382 第 10 :章边界层理论基础 量方程简化为 坠旦-1-坠卫 =0dX . dy pu 窍 +pv 去 =-2+  ($\mu$ 去) 学 = O+O( $\epsilon$2 ) cJ y 能量方程的简化过程中考虑到 UEEr u22. |当 1>> 1 旦旦 1. 1 旦旦 1>> 1 旦旦 1.

\[eJ y ,,-- dX' 1\theta  y 1" 1 d.r l' 1 dy 1" 1 a.r 1 '\]

|立 /$\mu$ 红、 1>> 1 主 /$\mu$ 江| dy 飞 $\theta$ yl $\theta$r 飞 dX 11 由此得到可压缩边界层内能量方程的简化形式为 $\iota$$\rho$uf+couf=ufE 十三 (${\rm A}$$\tilde{T}$)+u(\textasciitilde{}叫 Z.. o.r .. oy ox oy 飞 oy I 飞 (/y I (10.50) (10.5 1) (10.52) (10. 53a) 令 H 表示总婚，在边界层中有 :h=c ${\rm þ}$ T=H-u 2 /2 ， 将它代人式 (10. 53a) ，并与 u 乘动量方程所得动能方程相减，可得到边界层内总始形式的能量方程 z 或 pu 去 +$\rho$ 弩 = :y[$\mu$  (Z+(1 一生)三 )J (10. 53b) 边界条件 ${\rm 1}$速度边界条件

\begin{equation*}y=O:u=v=O; y \rightarrow \infty  : u=U.( .r)\tag{10.54}\end{equation*}

${\rm 2}$温度边界条件

\begin{gather*}
y=O:T=T ", =O (10. 55a)\\
aT (10. 55b) y=OZEY=qw
\end{gather*}

式 (10. 55a) 和式(1 0. 55b) 分别对应给定壁面温度或给定壁面热流 z 边界层外缘的温 度边界条件是

\[y \rightarrow \infty  : T= T.\]

边界层内的压强 $\rho$= 户. ,dp.ld.r= -p.U.dU./dx. 2. 可压缩流体定常层流边界层的主要特性 (1) Pr=l 的可压缩边界层能量方程的积分 (10. 55c) 可压缩边界层方程的动量方程与能量方程是藕合的，需要联立求解.在 Pr=l 的条件下，能量方程有以下的特解 z

T = T.(1 + 乓1叫)- U 2 00.56) 飞\textasciitilde{} 'C${\rm þ}$ 当 Pr=1 时，方程 00. 53b) 中 (1 -1/ Pr)u 2 /2=0 ，总熔形式的能量方程简化为

\[aH , aH a 1. aH \]

pu 百-; -r pv 百歹=巧$\tilde{A}$ iJy) 显然，此方程有一个特解 :H= 常数 z 而常数可由边界层外边界的衔接条件确定: 1 U! ., 1 u2 H = H. = h. + $\tau$=...!. =h+ $\tau$ 一 $\tilde{C}$ p Jt. Cp 或写成温度形式 z h = c${\rm þ}$T = c 几+号 -5 将 c${\rm þ}$T. 用声速表示(见第 7 章) ，就可得式(1 0.56). 在壁面上 u = 0 ，代人式 00.56) 得壁面温度 z TJtf=1(I+ 平时) 即 Pr=l 时壁面温度等于边界层外缘的绝热滞止温度.当 Pr =l=- l 时，绝热壁面的温 度和普朗特数有关，一般有 Tu. = T. + 去 = T. (1 + r( 丹)与1 险; ) r(Pr) $\zeta$1 称为恢复系数，就是说绝热壁面的温度小于边界层外缘的绝热滞止温度. (2) 可压缩速度边界层厚度的估计 根据不可压缩边界层厚度公式，可近似估计可压缩速度边界层厚度为 :8- x/(pu.x/$\mu$) 1/ 2 ,p ， 户为边界层内平均密度和平均粘度.因边界层内压强在法向为常 数，由状态方程可得 : p OC l/T oc l/h. 而且平均粘度可由幕次律计算: p=$\mu$.(T/T.)n C汇 p.(ii /h ， )n ， 因此 p/严\textasciitilde{}怡，巾， )(h ， /h)nll ， 代人厚度估计式，得 8c汇 x 1 主- \textbackslash\{\}号=一王一/豆\}哼l

\[(pU ,x /\mu ) I! Z   h. J /Re:   h. J\]

式中 Re... 为边界层特征雷诺数.当 Pr 等于 l 时，可由能量方程近似积分得到 h= h ， +U!/2- -;7 /2> 儿，因 U\textasciitilde{}/2-U2/2>0 ， 所以 h/h ， >I ， 于是有 $\delta$ $\infty$ x 1 h 飞平 、 ( !-) ---\_.一.、 mm dz lh，//飞 /Re: I;n 式中下标 com 表示可压缩， ln 表示不可压缩.容易理解马赫数越大，叨 /2- u2 /2 值 越大，再 /h ， 也越大，就是说当雷诺数不变时，高马赫数可压缩边界层厚度远大于不可 压缩边界层的厚度，即 8."m>>$\iota$. (3) 温度边界层厚度估计 在气体边界层中温度和速度一样从壁面渐近变化到来流温度，称为温度边界层.

可以通过能量方程的量级分析估计温度边界层. 若温度边界层厚度用 8T 表示，它的量级也是:品 /L=e<< l.将温度边界层方程无 量纲化，温度特征量为 : U\textasciitilde{} ， / 句，其他特征量和速度边界层相同.代入能量方程并整 理可得 pafL$\mu$ aT = 丘 \textasciitilde{}+J. \_${\rm A}$二二一立 (i iJ T \textbackslash\{\}\_j.\_ \_\_\_E.主E\_\_(但 f iJ.i ' t' $\tilde{a}$\textasciitilde{}豆 a .i ' eZ c${\rm þ}$p" Lu.. a y 飞.. a 豆 /ttpp Lu n 飞 3 豆/ 各项量级应当相当，最后两项的系数必须同一量级 z 1 ${\rm A}$二$\iota$一 =l\_\textasciitilde{}\textasciitilde{}=l\_\_!\_ 土 eZ C p$\rho$ " Lu " e2 p. Lu 咱 c p$\mu$ e 2 Re Pr 所以温度边界层的厚度应为 e2 cx: ($\delta$!:)Z cx:: \textasciitilde{} 1-::-'- I CX: '::一-一 飞 L I -- Re Pr 即温度边界层的厚度有以下估计: $\delta$r 民 LJ二二 而速度边界层厚度是: 8T cx: L/ .jR;.所以两种边界层厚度之比为 品./ 8 CX: 1/ ffr 由此可得出结论:温度边界层的厚度和速度边界层厚度之比与普朗特数的平方根成 反比.常见气体 : Pr<l. 因此 8T >8. 通过以上分析，高速气体的可压缩边界层具有以下特点. (1)在壁面上气体速度等于零，对边界层流动产生强烈的滞止作用，由于流体的 粘性耗散，物面附近的气体温度升高; (2) 可压缩气体速度边界层厚度大于相等雷诺数下的不可压缩边界层; (3) 气体绕流的温度边界层厚度大于速度边界层厚度; (4) 气体边界层内，流体的密度、粘度和导热系数都是变茧，这种情况下，连续方 程、动量方程和能量方程是锢合的，必须联立求解.

\section*{10.6  绕平板定常湍流边界层}\addcontentsline{toc}{section}{10.6 绕平板定常湍流边界层}

采用布西内斯克涡粘模型封闭雷诺应力，二维定常雷诺方程可写作(省略平均 符号) : u 争 +v\textasciitilde{}坐 =-l\_ \textasciitilde{}+主 12($\nu$ +v，) 争 1 dX dy pdX rJ.rL dXJ + [($\nu$+ 川(艺+去 )J (10. 57a)

u 去 +UZ=-j 去+ [($\nu$+ 川(艺+去 )J +主 12(11+$\nu$J f叮 dYL aYJ (1 0. 57b)

\[nu --h\]

一句+ u-z 、。一、。 (1 0. 57c\} 式中 $\nu$t 是涡粘系数.除贴近壁面外，灿》 $\nu$ ，一般情况，灿 /$\nu$-10 2 ${\rm •}$ 即湍流的涡粘系数总 是远大于分子粘性系数.在 Re，. =u.${\rm ‘}$ :X /II>> 1 的情况下，以涡粘系数为特征量的雷诺 数仍然很大，即仍有 (Reh= U..\_ L/ $\nu$ ， >>1. 例如当 Re l.= U山 :X /II= 108 .ReT -10 6 ${\rm •}$ 对于 湍流平均运动来说，仍然存在以湍流粘性主宰的很薄边界层.层流边界层运动的概 念，如边界层厚度等，在平均意义上可推广到湍流边界层中.利用边界层近似的方法， 忽略流向的层流和湍流粘性输运，可导出湍流薄层方程 z iJ u , dU 1 iJ tJ , d ${\rm i}$ / , , dU l 一 +$\nu$ 一=一$\tilde{vy}$ 十一 1 ($\nu$ 十川:\textasciitilde{} 1

\[\theta  :x ' - iJ y P J:x ' JyL - - -, JyJ\]

主 ($\rho$+< 山')\} = 0 rJy iJ u . a 守' :二 + :v = 0 J :X . dy (1 0. 58a\} (10.58b) 00. 58c\} 可以用混合长度模式或 h 模式计算涡粘系数，但是需要用数值计算方法得到 以上方程的解，它属于计算流体力学的内容.作为一种简易的近似方法，动量积分方 程是常用的计算湍流摩阻的丁.程方法.从雷诺平均的湍流边界层方程出发，可导出揣 流边界层的平均动量积分方程为(读者可仿效层流边界层的卡门-霍华斯方程的推 导，导出以下方程) :

\[d8z , 82 / n , T" dU,. < U'\]

\_,V< + ::: (2 + H) 一一=一一 (10.59a) d.l' ' U, ,- , ${\rm •}$ -, d.l' pU; 式中 82 为湍流边界层平均流动的动量厚度$\cdot $凡是壁面剪应力 ， H =81 /82 称为湍流边 界层的型参数. 1由流平板边界层流动中 U，=u.$\lambda$ = 常数，因此方程<1 0. 59a\} 可简化为

\[d8z <U'\]

d:x pU\textasciitilde{}.c (1 0. 59b\} 82 和 $\tau$E 分别为 品 =fd:二 (1- J. )dy (1 0.60) v 、${\rm ‘‘‘}$，，，，，u-y 3u - 、d

\[''''E\cdot \]

$\mu$ 一一u r 00.6 1) 式中 U(y) 是湍流边界层中平均流向速度分布， $\mu$ 是分子粘性系数.如果已知湍流边 界层中的平均速度分布，则由以上三式就可以求出边界层厚度和壁面剪应力的分布.

流平均运动的速度分布为对数型 z 旦 =Aln 旦王 +B U 营 $\nu$ (10.62) 式中 Ur = .. /rw/p 是壁面摩擦速度，在平板边界层中对数速度剖面表达式中常数 A.B 和圆管湍流中略有不同，常用值为 : A=O. 4.B=5. O. 将对数速度分布公式 00.62) 代人式 00.60) ，得 3=j;61l-6三 )d 去 =是I 二 (27 -t)d 去 =A 走 -2A 2 ( 在 f 另外，由对数速度分布可以得到 u .l U晤 的隐式关系，将 S 代人式 (10.62) .得 最后，式 (10.59b) 可改写为 旦旦 = Aln 但王 +B U 宫 $\nu$

\[dSz -ui\]

dx U\textasciitilde{}， (10.63) 00.64) 00.65) 方程(1 0.65) 的初始条件是 : X=0'\textasciitilde{}2=0. 联立式(1 0.63)- 式 00.65) 就可解出 三 个 未知量以 x) '\textasciitilde{}2 (川和 U.(x). 由于式 00.64) 是隐式代数方程，以上联立方程只能用 数值方法解出.在大量数值和实验结果的基础上，人们将式(1 0.64) 用近似显式公式 表示 z (走 ) E =o 00 叫半)

\[-'1/6\]

将它代入方程(1 0.65) 就可得到 \textasciitilde{}2 (x) 和战 (x) 的显式公式 z (半)= 0 。叫于 f l7 以及 (走 f = 0 川(平

\[)\nu 7\]

00.66) 00.67) 表面局部摩阻系数 Cf=-2一 =2 旦二，由式(1 0.67) 很容易导出局部摩回系数公式J pU'!.o /2 U乙 如下 z C乌$\leftarrow$川f卢$\lambda$=$\rightarrow$O 叫牛r

\[l\gamma -\rightarrow - 11\upsilon /\]

长为 L 的平板平均摩阻系数 C D 由式 (1 0.66) 积分求出 z 00.68)

\begin{gather*}
rL~. - -~-~/U\infty  L   -' 1/ 7\\
Co = iJo Cfdx = O. 0307 ~=~) (1 0.69)
\end{gather*}

湍流边界层中也可应用平均速度的幕画数分布公式.借用圆管湍流中速度分布 的事画数近似和壁面剪应力的布拉休斯公式 z 和 击= (去) 且可 =0.0225(Y$\tilde{J}$ $\rho$ U\textasciitilde{} 飞 U0 81 当 Re=UoR/ $\nu$<10 6 时 .71=1/7. 因为壁面附近的湍流平均运动有相似的性质，将以上 公式中矶和 R 分别换成边界层中的 U$\infty$ 和 8. 就得到平板湍流边界层速度分布的罪 函数近似式如下: Joo = (去 r 以及 iL=0.0225i-LI pU\textasciitilde{} 飞 Um 81 将速度分布公式代人动量厚度公式 00.63) 得

\[82 = 78/72\]

将它和剪应力公式代人动量积分式 00.6 日，得

\[1 d8   ~ -\wedge \wedge ~ 1 \nu    1/4\]

(7/72) ( -;v 1 = 0.0225 (一$\tilde{l}$ \textbackslash\{\}dx) 飞 U$\infty$ 81 积分的初始条件为: x=0.8=0. 积分结果得 8 = O. 37($\iota$11 生\textasciitilde{}) 4/5 飞 U oo I 飞 $\nu$/ 代人剪应力公式，得局部摩阻系数 Cr= \_:，w.\_=O …。/豆豆王 $\gamma$ 1 /5一一…..\textasciitilde{} . - . f-FDE-hud ，口 '-v-j 长 L 平板的平均阻力系数 I CrdX .rT Y - 11, In '-"f V.J. \_ \_\_\_IU 二 L \textbackslash\{\} - 1/.

\[C/) =\iota \tau - = 0.072(~\iota l\]

实验证实，在 Rex <107 时，速度分布的幕函数近似公式是足够好的 z 当 Rex >10 7 后，对数速度分布公式和实际符合较好. 和层流边界层对比，湍流边界层厚度沿流向的增长为 8.8$\infty$ (v/U，$\infty$ ) (Re;r) 俏，它远 远大于层流边界层 8. 8OC ($\nu$/U\textasciitilde{} ) (Re;r )1 /2 ; 湍流边界层的壁面摩擦系数 Co(Cooc(U，$\infty$ L/ $\nu$) - 1/ 5 )也远远大于层流边界层的摩擦系数 Co(Cooc(U$\infty$ L/$\nu$) - 1 / 2 ). 读者可以发现应用动量积分公式计算捕流边界层时和求解层流边界层的动量方

程有所不同，我们除了给出近似速度分布外，还需要有壁面摩擦系数和湍流特征量之 间的关系.这是因为湍流边界层中的'Z"" 等于等雷诺应力层中的雷诺应力.它需要模 型来封闭.在应用动量积分方程求解时.就需要补充壁面摩擦应力的经验关系式. 最后，在实际边界层流动中，从边界层前缘开始.流动是由层流发展到湍流.因此 计算时，需要判断何处开始边界层是湍流.前面已经陈述过.流动从层流到湍流是 - - 个过程，而这一过程十分复杂，至今尚无法用理论方法来预测.常用的方法是线性稳 定性理论和实验相结合来提供层流到湍流转披点的经验性的预估公式，井假定转缺 点以前边界层完全是层流，转换点以后边界层完全是湍流.从流动的线性稳定性理论 的讨论中，我们已经知道，平行流动的失稳主要和流动的速度剖面以及特征雷诺数有 关，在平板边界层中上述两个因素都和边界层的厚度有关. 一 个丁.程常用的转缺点的 预估公式为塞贝阿史密斯 (Cebeci - Smith) 公式: rJ $\rho$ f ， , 22400 、 (Re o)" = 一二二=1. 174(1+ 一一一'，" j(Re ,):';" ; (10.70) $\nu$ 飞 (Re，)" J 式中 Re. =U,.r / $\nu$ 是边界层的流向需诺数 .0 是边界层的动世损失厚度，下标川 r" 者，民 转缺点.式(1 0.70) 的具体应用如下=从层流边界层计算开始，计算每个流向位置 I 的动量损失厚度。(川和相应的雷诺数 Reo . 当 Red 和 Re ， 满足式(1 0.70) 时，就认为民 流边界层转换为湍流.在转缺点后 .1" > 儿，应用湍流边界层方程进行计算.

\section*{习题}\addcontentsline{toc}{section}{习题}

\noindent\textbf{10.1.} 证明:边界层内任意 一 个.r = const. 的截面上有

\begin{gather*}
(1 )\delta >o,  > \delta :? ;\\
(2) o>o , + \delta 2 ;
\end{gather*}

上述结论是杏适用于不定常流动?是否适用于可压缩流动?

\subsection*{例 0.2  μ=0. 731Pa. s.p=925kg / 旷的油，以 0.6m / 目的速度平行地流过 一 块}

氏为 0.5m 、宽为 0.15m 的光滑平板.求边界层最大厚度及平板所受的阻力.

\noindent\textbf{10.3.} 假定层流边界层的速度分布如下，请计算边界层的排挤厚度 ${\rm o}$ ，和l 动!d:损

失厚度 $\delta$2 . (1) u=U\textasciitilde{} , y<仇 u=U ， y > ${\rm o}$; (2) u=ul 三旦 \_ \textasciitilde{} ( \textasciitilde{} )3 l, v <$\delta$ ; u= U , v>${\rm o}$. L 2 ${\rm o}$ 2 飞 ${\rm o}$ J J' J \textasciitilde{} \textasciitilde{}， - -. J

\noindent\textbf{10.4.} 用动量积分关系式计算 10.3 题假定的边界层速度分布的壁面摩阻系数.

\noindent\textbf{10.5.} 假设平极层流边界层内速度分布为 u / U， =u + bsin( cy / 的，式中 u.b.c 为

待定常数，试用动量职分关系式方法求 ${\rm o}$， / ${\rm O}$2 和 C/J.

\section*{思考题}\addcontentsline{toc}{section}{思考题}

\noindent\textbf{10.6.} 已知层流边界层的外边界速度分布队 =u，.\_ (1- .r/ U. 求分离点位置

I，/ L. 又若流体的运动粘性系数为 $\mu$ ，求 I /L =0.0888 处的 81 .82 和 C ，.

\noindent\textbf{10.7.} 假设平报湍流边界层内统计平均速度分布为1/ 10 次方规律，并设全平极

为湍流，用动量积分关系式证明其结果为: 8/:J\_' = o. 239Re.. 0. 15.1 ${\rm •}$ C" =0. 0362Re ,. o. 101 ${\rm •}$

\noindent\textbf{10.8.} 一 块长为 6m 、宽为 2m 的光滑平板，平行地放置在来流速度为 60m / s. 温

度为 293K 的空气流中，已知边界层流动的临界 Re 为 10 6 (即 Re<10 6 用层流公式计 算 .Re > 10 6 用湍流公式计算 ).求平极所受阻力.

\noindent\textbf{10.9.} 温度为 293K 的空气以 30m/s 速度平行地流过 2m 长的光滑平板，已知

临界 Re 为 5 X 10 5 ${\rm •}$ 试估算离平板前缘 1m 处的边界层厚度以及尾缘处的边界层 厚度.

\noindent\textbf{10.10.} 流线型火车长 110m 、宽 2.75m 、商 2.75m. 假定顶部及所布侧面的表面

摩擦力相当于长 110m 、宽 8.25m 的光滑平极的 一 个表面.火车以 160km / h 的速度 在空气巾等速行驶，唁气密度 p = 1. 22kg/m 门粘度 $\mu$= 1. 79 X 10 - f. Pa ${\rm •}$ S. 如临界 Re 为 5 X 10 ' . 问层流边界层能维持多长?火丰为克服风的摩擦阻力.消耗的功率为 多少?

\subsection*{例 0.11  10.10 题中如选用1/ 10 次方速度分布规律，试计算火车为克服风的摩}

擦阻力所消柜的功率'. S10. 1 在风洞 t 作段(没有模型)巾要达到流向 lli 强梯度等于零.可以采用外扩 工作段壁面的方法，为什么?如何计算扩张最?也可采用壁面吸气的方法，为什么? 吸气量怎么计算? S10.2 温度边界层的厚度和速度边界层厚度之比，哪个大? S10.3 在高雷诺数回球绕流中，为什么提高来流的湍流度可以使阿球阻力减 小?为什么在前半网球处施加壁面扰动，也可使图球阻力减小? S10.4 在机翼绕流或其他钝体绕流中，分离点附近吸气或吹气都能减小阻力. 为什么?
